%%%%%ENTREES%%%%
\begin{tabarg}
\multicolumn{3}{l}{\titretab {Entr\'ees}}\\
\hline
$X$ & Vecteur, \texttt{float} & Param�tre � s�lectionner (unit� quelconque)\\
$Y$ & Vecteur, \texttt{float} & Param�tre selon lequel s'effectue la s�lection (unit� quelconque)\\
$Crit\_{min}$ & Coeff., \texttt{float} & valeur mini de $Y$ pour laquelle $X$ est consid�r� valide\\
$Crit\_{max}$ & Coeff., \texttt{float} & valeur maxi de $Y$ pour laquelle $X$ est consid�r� valide\\
$Y_{inval}$ & Coeff., \texttt{float} & valeur de $Y$ qui permettra d'�liminer �galement des valeurs de $X$ (�galement pour le cas o� la "valeur manquante" est mal reconnue (attribut misval non renseign� dans $Y$)) \\
$Inv\_val$ & Coeff., \texttt{float} & valeur de remplacement de $X$ lorsque $Y$ ne correspond pas aux crit�res, le plus souvent "valeur manquante", 
soit \texttt{\$Arg1.misval}\\
\end{tabarg}

\medskip
%%%%%SORTIES%%%%
\begin{tabarg}
\multicolumn{3}{l}{\titretab {Sortie}}\\
\hline
$X_{valid}$ & Vecteur, \texttt{float} & $X$ lorsque $Y$ correspond aux crit�res. $Inv\_val$ dans les autres cas. 
(m�me unit� que $X$).\\
\end{tabarg}
